\documentclass{article}

\PassOptionsToPackage{numbers,compress}{natbib}
\usepackage{neurips_2026}

\usepackage[utf8]{inputenc}
\usepackage[T1]{fontenc}
\usepackage{hyperref}
\usepackage{url}
\usepackage{booktabs}
\usepackage{amsfonts}
\usepackage{nicefrac}
\usepackage{microtype}
\usepackage{xcolor}
\usepackage{amsmath}
\usepackage{graphicx}
\usepackage{amssymb}
\usepackage{enumitem}

\title{Post-Training is About States, Not Tokens:\\
A State Distribution View of SFT, RL, and On-Policy Distillation}

\author{%
  Anonymous Authors\\
  Affiliation\\
  Address\\
  \texttt{email} \\
}

\begin{document}

\maketitle

\begin{abstract}
Large language model post-training methods such as supervised fine-tuning (SFT),
reinforcement learning (RL), and distillation are often analyzed through their
loss functions: maximum likelihood, policy gradients, forward KL, reverse KL, or
related objective-level variants. We study a complementary factor: the
\emph{state distribution} on which supervision is applied. For an autoregressive
policy, a state is a prompt plus generated prefix. SFT trains on fixed dataset
states, while RL and on-policy distillation (OPD) train on states induced by the
current learner.

We formalize post-training as state-distribution shaping and run a controlled
small-scale study using Qwen3-0.6B-Base on GSM8K, with TruthfulQA and MMLU as
retention evaluations. Our results show three phenomena. First, a mild SFT run
improves GSM8K with little forgetting, while a stress SFT run causes substantial
retention loss. Second, OPD from a degraded SFT teacher surpasses that teacher
on GSM8K, TruthfulQA, and MMLU, despite using the teacher as its only
supervision source. Third, a lightweight on-policy RL run improves GSM8K while
preserving retention. These results support a state-centric view of
post-training: the source and locality of training states can be as important as
the form of the supervision signal.
\end{abstract}

\section{Introduction}

Post-training is central to modern language model development. Supervised
fine-tuning teaches models to imitate demonstrations, reinforcement learning
optimizes model generations against rewards, and distillation transfers behavior
from a teacher to a student. These methods are usually compared by their
objectives: maximum likelihood for SFT, reward maximization for RL, and
token-level divergences for distillation.

This view is useful, but incomplete. In autoregressive generation, each update is
not merely a token-level update. It is an update at a \emph{state}: the prompt
and the prefix that the model has already generated. The distribution over these
states determines which parts of the policy are touched during training. SFT
updates the model on dataset prefixes, which need not be states the current
policy would visit. RL updates the model on states sampled from the current
policy. OPD occupies a middle position: the teacher provides the signal, but the
student determines the states.

We investigate the following thesis:
\begin{quote}
Post-training behavior is governed by the interaction between supervision signal
and training state distribution; on-policy state supervision can preserve
locality and allow students to improve beyond degraded teachers.
\end{quote}

We make three contributions:
\begin{enumerate}[leftmargin=*]
    \item We formulate SFT, RL, and OPD under a common state-distribution view of
    autoregressive post-training.
    \item We implement a controlled single-GPU experimental pipeline measuring
    target accuracy, retention, forgetting, and rollout-state drift.
    \item We provide evidence that OPD can outperform a degraded SFT teacher and
    that on-policy RL improves target performance with little retention loss.
\end{enumerate}

\section{State Distribution View}

\subsection{Autoregressive States}

An autoregressive language model defines a policy
\begin{equation}
    \pi_\theta(y_t \mid x, y_{<t}).
\end{equation}
We call
\begin{equation}
    s_t = (x, y_{<t})
\end{equation}
the \emph{state}. The next token $y_t$ is the action, and a generated answer is a
trajectory through states. Let $d^\pi(s)$ denote the state visitation
distribution induced by policy $\pi$ on a prompt distribution.

\subsection{SFT: Off-Policy State Fitting}

SFT minimizes token loss on dataset states:
\begin{equation}
    \mathcal{L}_{\mathrm{SFT}}(\theta)
    =
    -\mathbb{E}_{s \sim d_{\mathrm{data}}}
    \log \pi_\theta(y^\star \mid s).
\end{equation}
The supervision is dense, but the states are off-policy. If dataset trajectories
are far from the model's own rollouts, updates may affect behavior in regions
that the model cannot reliably reach or recover from. This is the familiar
exposure-bias problem recast as a state-distribution mismatch.

\subsection{RL: On-Policy Local Improvement}

RL samples trajectories from the current policy and updates the model using
rewards:
\begin{equation}
    \max_\theta \mathbb{E}_{s \sim d^{\pi_\theta}, y \sim \pi_\theta(\cdot|s)}
    [r(s,y)].
\end{equation}
The reward may be sparse, but the states are on-policy. This makes RL a local
improvement procedure: it modifies behavior where the current model actually
visits.

\subsection{OPD: Teacher-Guided On-Policy Learning}

In OPD, the student samples states and the teacher provides supervision:
\begin{equation}
    \mathcal{L}_{\mathrm{OPD}}(\theta)
    =
    \mathbb{E}_{s \sim d^{\pi_S}}
    \left[
    D(\pi_T(\cdot|s) \,\|\, \pi_S(\cdot|s))
    \right].
\end{equation}
In our strongest OPD variant, the teacher generates short continuations from
student states and the student learns those continuations with cross entropy.
This is analogous to DAgger-style learning: the learner controls the state
distribution, while an expert-like source provides local repair signals
\citep{ross2011dagger}.

\section{Unified Framework}

We view post-training as repeatedly transforming a model's state distribution:
\begin{equation}
    d_{k+1}(s) = \mathcal{T}\big(d_k(s), \mathrm{signal}\big).
\end{equation}
Different algorithms vary both in the source of states and the source of signal.

\begin{table}[t]
\centering
\caption{State-source view of common post-training methods.}
\label{tab:framework}
\begin{tabular}{lll}
\toprule
Method & Training state source & Supervision signal \\
\midrule
SFT & Dataset trajectories & Gold tokens \\
Offline KD & Teacher trajectories & Teacher logits/tokens \\
OPD & Student trajectories & Teacher logits/continuations \\
RL & Current policy trajectories & Reward \\
DAgger & Learner trajectories & Expert actions \\
\bottomrule
\end{tabular}
\end{table}

This view predicts that methods using learner-induced states can behave
differently from off-policy imitation even when the supervision source is similar
or weaker. In particular, a student can outperform its teacher if the teacher's
errors are coupled to the teacher's own state distribution rather than fully
encoded in its local responses on the student's states.

\section{Experimental Setup}

\paragraph{Model and hardware.}
All experiments use Qwen3-0.6B-Base with LoRA adapters on a single RTX 3090
24GB GPU. We use plain GSM8K-style prompts rather than chat templates because
the model is a base model.

\paragraph{Target and retention tasks.}
The target task is GSM8K. Retention is measured with TruthfulQA multiple-choice
and a selected MMLU subset. Base model scores are GSM8K 0.448, TruthfulQA 0.300,
and MMLU 0.436.

\paragraph{Methods.}
We evaluate mild SFT, stress SFT, OPD variants, and lightweight on-policy RL. The
mild SFT run uses GSM8K SFT data and produces a non-degraded teacher. The stress
SFT run uses five epochs, learning rate $5\times10^{-4}$, LoRA rank 64, and LoRA
alpha 128, intentionally probing forgetting. OPD samples student states and
trains from teacher continuations. RL uses group-relative exact-answer reward on
GSM8K rollouts.

\paragraph{Metrics.}
For a retention task, forgetting is
\begin{equation}
    F = \mathrm{Score}_{\mathrm{base}} - \mathrm{Score}_{\mathrm{post}}.
\end{equation}
Retention ratio is $\mathrm{Score}_{\mathrm{post}} /
\mathrm{Score}_{\mathrm{base}}$. We report mean forgetting and mean retention
over TruthfulQA and MMLU.

\paragraph{State drift.}
For each trained model, we sample rollouts on a fixed prompt set and collect
prefix states $s_t=(x,y_{<t})$. We embed states with a lightweight lexical
feature representation and report maximum mean discrepancy (MMD) with an RBF
kernel. We also compute centroid distance, sliced Wasserstein distance, and
lexical Jaccard distance, but use MMD as the primary scalar drift measure.

\section{Results}

\subsection{SFT Can Be Gentle or Destructive}

\begin{table}[t]
\centering
\caption{Main results. GSM8K is the target task; TruthfulQA and MMLU are
retention tasks. Forgetting and retention are averaged over TruthfulQA and
MMLU.}
\label{tab:main-results}
\begin{tabular}{lrrrrrr}
\toprule
Run & GSM8K & TruthfulQA & MMLU & MMD & Forgetting & Retention \\
\midrule
Base & 0.448 & 0.300 & 0.436 & -- & -- & -- \\
Mild SFT & 0.512 & 0.295 & 0.444 & 0.00956 & -0.0015 & 1.0008 \\
Stress SFT & 0.420 & 0.245 & 0.364 & 0.01093 & 0.0635 & 0.8258 \\
OPD from mild SFT & 0.512 & 0.290 & 0.434 & 0.01470 & 0.0060 & 0.9810 \\
OPD from stress SFT & 0.466 & 0.275 & 0.430 & 0.01092 & 0.0155 & 0.9515 \\
On-policy RL & 0.472 & 0.290 & 0.442 & 0.01098 & 0.0020 & 0.9902 \\
\bottomrule
\end{tabular}
\end{table}

Mild SFT improves GSM8K from 0.448 to 0.512 with essentially no retention loss.
This is an important negative control: SFT does not necessarily forget in our
setup. However, stress SFT produces substantial retention degradation:
TruthfulQA falls from 0.300 to 0.245 and MMLU falls from 0.436 to 0.364. Its
mean retention ratio is 0.8258. Interestingly, stress SFT also reduces GSM8K to
0.420, indicating that overly aggressive off-policy training can degrade both
general and target behavior.

\subsection{OPD Can Surpass a Degraded Teacher}

The clearest OPD result uses the stress SFT model as the teacher. The teacher is
degraded: GSM8K 0.420, TruthfulQA 0.245, MMLU 0.364. OPD from this teacher
achieves GSM8K 0.466, TruthfulQA 0.275, and MMLU 0.430. Thus the student
surpasses its teacher on all measured tasks despite using that teacher as its
supervision source.

This supports the claim that teacher behavior is not transferred as a single
global object. The student receives local guidance on states sampled from the
student's own policy. When those states differ from the teacher's problematic
trajectory distribution, the student can avoid inheriting some teacher failures.

\subsection{RL Provides the On-Policy Reward Point}

The lightweight RL run improves GSM8K from 0.448 to 0.472 while retaining
TruthfulQA 0.290 and MMLU 0.442. Its mean forgetting is only 0.0020. This is
consistent with the view that RL behaves as an on-policy local improvement
method: it changes behavior where the policy samples states and does not require
large off-policy movement.

\subsection{Drift Magnitude Is Not the Whole Story}

Our initial hypothesis was that scalar state drift would strongly explain
forgetting. The data suggest a more nuanced conclusion. Stress SFT and OPD from
the stress teacher have nearly identical MMD drift, 0.01093 and 0.01092, but very
different retention ratios, 0.8258 and 0.9515. Similarly, RL has comparable MMD
drift, 0.01098, with much smaller forgetting.

Therefore, the evidence does not support a simple scalar law of the form
``larger MMD implies more forgetting'' in this small setup. Instead, it supports
a state-source claim: the \emph{quality, locality, and learner-dependence} of
training states matter. Measuring only the distance between rollout
distributions can miss whether updates were applied on states that are locally
recoverable for the learner.

\section{Discussion}

\paragraph{Objective-level analysis is incomplete.}
The contrast between stress SFT and OPD from the stress teacher is difficult to
explain using only the source of supervision. Both are ultimately shaped by the
same degraded teacher/data behavior, but OPD applies supervision on states
sampled from the student. This changes the learning problem.

\paragraph{On-policy dense shaping.}
The successful OPD variant used teacher continuations rather than one-step
logit matching. One-step OPD collapsed badly in our runs, reaching GSM8K 0.040.
Continuation-based OPD recovered target performance. This suggests a practical
recipe: combine on-policy sampling with dense, trajectory-level local
supervision.

\paragraph{Limitations.}
This study is small-scale: one base model, one target dataset, LoRA adapters,
limited retention tasks, and lightweight drift estimators. The RL trainer is a
minimal on-policy GRPO-style implementation rather than a full-scale verl PPO or
GRPO setup. Our drift metric uses lexical features rather than hidden-state or
encoder embeddings. The results should therefore be read as evidence for a
mechanistic hypothesis, not as a benchmark claim.

\section{Related Work}

\paragraph{Post-training objectives.}
Language model post-training is commonly described through the optimization
objective being used. SFT is typically framed as maximum-likelihood learning on
instruction demonstrations. RLHF and related approaches optimize model samples
against learned or verifiable rewards \citep{ziegler2019fine,ouyang2022training}.
Preference-optimization methods such as DPO remove the explicit online RL loop
and express preference learning as a supervised objective
\citep{rafailov2023direct}. This objective-centric view has clarified many
algorithmic trade-offs, but it can obscure the role of the state distribution on
which the objective is applied. Our work keeps the objective visible, but treats
the training state source as a separate axis.

\paragraph{Exposure bias and imitation learning.}
The distinction between dataset states and learner-induced states has a long
history in imitation learning. Behavioral cloning trains on expert trajectories
and can suffer from compounding errors when the learner visits states absent
from the demonstrations. DAgger addresses this by collecting learner-induced
states and querying an expert on those states \citep{ross2011dagger}. Exposure
bias in sequence prediction captures a related mismatch between training on
gold prefixes and testing on model-generated prefixes
\citep{bengio2015scheduled}. We adapt this perspective to LLM post-training:
SFT resembles behavioral cloning on fixed trajectories, whereas RL and OPD apply
signals on states induced by the current learner.

\paragraph{Knowledge distillation.}
Knowledge distillation transfers behavior from a teacher to a student through
soft targets, logits, or generated data \citep{hinton2015distilling}. In
language models, distillation is often analyzed by the divergence being
optimized or by the quality of teacher outputs. Our OPD experiments isolate a
different factor: the teacher may provide the signal, but the student can
control the state distribution. This explains why a student need not inherit all
failures of a degraded teacher, especially when teacher supervision is queried
on student states rather than copied from teacher trajectories.

\paragraph{Catastrophic forgetting and retention.}
Catastrophic forgetting has been studied in continual learning as the loss of
previous capabilities when training on new tasks \citep{kirkpatrick2017overcoming}.
In LLM post-training, forgetting is often measured as degradation on broad
retention tasks after specialization. Our experiments follow this empirical
tradition by measuring TruthfulQA and MMLU after GSM8K post-training. The
results suggest that forgetting is not only a matter of update magnitude or
dataset size: aggressive off-policy SFT can damage retention, while on-policy RL
and OPD can preserve more capability under comparable target-task pressure.

\paragraph{State drift measurement.}
Distribution shift is often quantified using embedding distances, classifier
two-sample tests, MMD, or Wasserstein-style metrics. We use rollout-state MMD as
a compact proxy for state drift, while also tracking other lexical distribution
statistics. Our results show both the value and limits of such scalar measures:
stress SFT and OPD from the stress teacher have nearly identical MMD drift, but
very different retention. This motivates richer state-distribution analyses that
consider not only how far model rollouts move, but which training states receive
supervision and whether they are locally reachable by the learner.

\section{Conclusion}

We proposed a state-distribution view of post-training. In this view, SFT is
off-policy state fitting, RL is on-policy reward-guided improvement, and OPD is
teacher-guided on-policy learning. Our experiments show that OPD can surpass a
degraded SFT teacher and that on-policy RL improves a target task with little
retention loss. The evidence refines the original thesis: post-training is not
only about token objectives or scalar distribution drift; it is about where
supervision is applied in the model's state space.

\bibliographystyle{plainnat}
\bibliography{references}

\end{document}
